\documentclass[12pt, a4paper]{scrartcl}
\usepackage[lithuanian]{babel}
\setlength{\marginparwidth}{2cm}
\usepackage{../../1.\space Vienanariai\space ir\space daugianariai/manostilius_lt}
\renewcommand{\theproblem}{\arabic{problem}}
\usepackage{geometry}
\usepackage{tasks}
\settasks{label={(\arabic*)}, label-width=2em, item-indent=3em}

\geometry{top=2cm, bottom=2cm, left=2cm, right=2cm}

\title{\textbf{Namų darbas: Antikos dvasia}}
\author{Andrius Berniukevičius}
\date{2026-10-03}

\begin{document}

\maketitle

\section*{Namų darbo atlikimo tvarka}

\begin{itemize}[leftmargin=*]
    \item Užduotis atlikite A4 formato lapuose ir atneškite penktadienį. Darbai bus surinkti.
    \item Sprendimus rašykite įskaitomai ir tvarkingai. Neįskaitomas, netvarkingas arba pribraukytas darbas nebus priimtas – jį reikės perrašyti ir tą pačią dieną atnešti.
    \item Atsakymus pažymėkite spalva arba apibrėžkite stačiakampiu.
    \item Jei užstrigsite, klauskite per „Messenger“ arba „Teams“ – padėsiu.
    \item Nepalikite darbo paskutinei dienai.
\end{itemize}

\section{Sąlygos}

\begin{problem}
Senovės Spartoje naudotis lanku buvo laikoma didžiausia gėda – tai bailių
ginklas, vengiantis artimos kovos. Tikras spartietis su priešu visada grūmėsi
akis į akį, pasikliaudamas tik savo kardu, ietimi, skydu ir drąsa.

Panašus yra ir licėjisto garbės kodeksas: atakuoti
sudėtingus matematinius uždavinius pasitelkiant „smegenų protezus“ (skaičiuotuvus, gpt, gemini)
ar buką jėgą (varginančią daugybą stulpeliu) yra žemas poelgis. Tikras
matematikos karys į kovą stoja apsiginklavęs vien algebros formulėmis, aštriu protu ir geležine logika.

Įrodykite, kad šio reiškinio reikšmė yra lygus legendinių Spartos karių
skaičiui, kurie Termopilų tarpukalnėje stojo į  kovą prieš
daugiatūkstantinę persų armiją:
\[
\frac{246^2+108\cdot246+(3\sqrt{6})^4}{277^3+23^3}
\cdot
\frac{277^2\cdot254^2-529^2}{277\cdot254-23^2}
\]
\end{problem}

\begin{problem}
Jaunas graikas atėjo į Delfų šventyklą, norėdamas tapti orakulo mokiniu.
Žynė jam padavė pergamentą su uždaviniu ir pasakė: „Išspręsk tai
naudodamas tik algebros dėsnius. Galutinis atsakymas atskleis, kas yra
pasaulio pradžia.“

Uždavinys atrodė taip:
\[
\frac{
(\sqrt{2}+\sqrt{3}-\sqrt{5}-\sqrt{7})
\cdot
(\sqrt{7}+\sqrt{5}+\sqrt{3}+\sqrt{2})
+\text{Olimpo dievų skaičius}+\sqrt{?}
}{
\dfrac{\sqrt{3}+\sqrt{2}}{\sqrt{3}-\sqrt{2}}
}
\]

Pamatęs klaustuką, jaunuolis paklausė: „Koks skaičius slepiasi po šaknimi?“

Orakulas atsakė:
\begin{itemize}
    \item „Esu trijų veidų skaičius, niekada neperžengiantis antrojo
    šimtmečio slenksčio.“
    \item „Sudėjęs visus mano veidus, gausi lygiai tiek, kiek yra tobulų
    Platono kūnų.“
    \item „Esu visiškai paklusnus septynioms Apolono lyros stygoms –
    padalintas iš jų, nepalieku jokio likučio.“
\end{itemize}

Įrodykite, jog šio reiškinio reikšmė nėra skaičius...
\end{problem}

\begin{problem}
Išskaidykite dauginamaisiais:
\begin{tasks}(2)
    \task $-5x ^{2} +4x +1$ 
    \task $x ^{3} + 3x ^{2} - 4x - 12$
    \task $x ^{3} +2x ^{2} -8x$
    \task $x ^{3}-16 x  + (3x-12)(x+1)$   
\end{tasks}
\end{problem}

\begin{problem}
    Žinoma, jog  daugianarį $x^3+4 x^2-36 x-144$ galima užrašyti kaip trijų pirmo laipsnio daugianarių sandaugą. Raskite tuos pirmos laipsnio dauginarius, jeigu vieno iš jų šaknis lygi $-6$.
\end{problem}

\begin{problem}
    Raskite kvadratinio trinario $x ^{2} +6x+c$ koeficiento $c$ reikšmę, jeigu jo šaknų kvadratų suma lygi $26$.   
\end{problem}

\begin{problem}
Duotas kvadratinis trinaris $ax ^{2} +bx+c$. Žinoma, jog $a$ sutinka su $c$ kaip $2 : 5$, o šių koeficientų skirtumas lygus $9$. Nustatykite koeficiento $b$ reikšmę, jei šio daugianario šaknys $x_1$ ir $x_2$ tenkina lygybę:
$$ \frac{1}{x_1} + \frac{1}{x_2} = \frac75.  $$           
\end{problem}

\begin{problem}
Išspręskite lygtis realiųjų skaičių aibėje:
\begin{tasks}(1)
    \task $\left(x^2-3 x+1\right)^2=4(x-2)(x-1)+1$
    \task $(3x-2)^2=(x^2+x-10)^2$
    \task $(x^2+x-10)^2+(6x-4)(x^2+x-10)+9x ^{2} +4=12x$  
    \task $x^5-2x^4+x^3-2x^2+x-2 = 0$ 
    \task $x^5+x^4+5x^3+8x^2+8x+40=0$
    \task $4^x = 16-6\cdot 2^x$
\end{tasks}
\end{problem}

\begin{problem}
        Kvadratinės lygties $ax ^{2} +bx+c = 0$ realiųjų sprendinių $x_1$ ir $x_2$ skirtumas lygus $5$, o kvadratų suma lygi $\frac{57}{2}$. Raskite $a+b+c$, jei $b>0$, o $a$ ir $c$ yra natūralieji vienaženkliai skaičiai.  
\end{problem}

\begin{problem}
Įrodykite, jog daugianaris $2x ^{2} +3x+10$ visuomet įgis tik teigiamas reikšmes.  
\end{problem}

\begin{problem}
    Su kuriomis $m$ reikšmėmis daugianaris $P(x) =  \frac{1}{4} x ^{2} +3x+2m ^{2} +m-1$ turės dvi vienodas realiąsias šaknis.  
\end{problem}



\newpage
\section{Atsakymai}

\begin{enumerate}[label=\textbf{\arabic* užduotis.}, leftmargin=*]
    \item \mbox{}
    \item \mbox{}
    \item
    \begin{enumerate}[label={(\arabic*)}]
        \item $(5x+1)(1-x)$
        \item $(x+3)(x-2)(x+2)$
        \item $x(x+4)(x-2)$
        \item $(x-4)(x^2+7x+3)$
    \end{enumerate}
    \item $(x+4)(x-6)(x+6)$
    \item $c=5$
    \item $b=-21$
    \item
    \begin{enumerate}[label={(\arabic*)}]
        \item $x_1=-1,\quad x_2=1,\quad x_3=2,\quad x_4=4$
        \item $x_1=-6,\quad x_2=-2,\quad x_3=2,\quad x_4=4$
        \item $x_1=-6,\quad x_2=2$
        \item $x_1=2$
        \item $x_1=-2$
        \item $x_1=1$
    \end{enumerate}
    \item $a+b+c=11+16\sqrt2$
    \item \mbox{}
    \item $m=2$ arba $m=-\frac52$
\end{enumerate}

\end{document}
